
Short birth spacing is more prevalent among disadvantaged social groups in India. \cite{singh2024wealth} finds that women from Scheduled Castes (OR = 1.18), Scheduled Tribes (OR = 1.14), and Other Backward Classes (OR = 1.12) have significantly higher odds of short birth spacing compared with forward-caste women. This may be due to well documented social group inequalities in coverage and awareness of family planning and maternal healthcare services, even after controlling for economic status. \citep{kumar2016explaining}. \cite{singh2024wealth} also finds that wealth inequalities were pronounced: mothers in the poorest quintile had nearly twice the odds of short birth intervals compared to those in the richest quintile. This raises the question of whether disadvantaged social groups in India are more likely to experience maternal depletion due to low wealth or other structural factors, and whether this can account for social group inequalities in prepregnancy health outcomes.

\begin{itemize}
    \item \cite{kuppusamy2023high}: women with no formal education and in the poorest wealth quintile were much more likely to have high risk pregnancy (NFHS-5)

    \item \cite{doke2022prevalence} found that among women desiring pregnancy in Nashik, selected risk factors were significantly more prevalent in tribal areas, where 7.7\% of high-parity women (4 or more births) still wished to conceive, compared to 3\% in non-tribal areas.
    
\end{itemize}
Additionally, maternal depletion is likely to be more severe for women who begin pregnancy with weaker nutritional reserves. In other words, the effects of low birth spacing and repeated pregnancies may be more severe for women of lower socioeconomic status are already at risk of underweight for other reasons. This interaction underscores why it is difficult to empirically isolate maternal depletion from socioeconomic disadvantage, as the two processes reinforce one another.

The gap between pre-pregnancy underweight of Muslim women that weren't SC or ST and forwards caste women was not statistically significant. Despite multidimensional disadvantage, Muslims have lower rates of pre-pregnancy underweight than forward caste Hindus, and this gap widens after accounting for wealth differences.


The hierarchal ranking of the caste system is based on ritual notions of purity and pollution ascribed to people, occupations, objects, and bodily processes. For example, Valmiki communities across North India have historically been confined to sanitation work, including manual scavenging, and regarded as ritually untouchable, subject to severe social exclusion \citep{valmiki2008joothan}. 



According to a 2021 Pew Research Center survey, 35\% of Indians self-identify as OBC, 25\% as Scheduled Caste, and 9\% as Scheduled Tribe \citep{starr2021measuring}. The same survey finds that the remaining share identify as belonging to the “General Category,” often referred to as forward castes. In this paper, we refer to individuals of Scheduled Castes as Dalits and individuals of Scheduled Tribes as Adivasis, using widely accepted terms of self-identification. Otherwise, we use the Indian Constitution’s administrative labels for consistency with demographic and administrative data, with no normative implications.


Because caste is the defining basis of status and power among the Hindu majority in India, class hierarchies among Muslims followed through conversion and assimilation. 


The dynamics between these social groups continue to be governed by notions of ritual purity, under which objects and people are considered polluting or purifying. For example, Balais in central India (a Dalit community traditionally associated with weaving) were not allowed to use village wells for water, under fear that they would spread impurity to other villagers \citep{ambedkar2014annihilation}.

\begin{itemize}

    \item \cite{gupta2022large}: SC and ST individuals have life expectancies between 4.2-7 years lower than high-caste individuals between 2013-2016. Muslim life expectancy was up to 2.8 years lower than high-caste in 2013-2016.

    \item \cite{gupta2020caste}: prevalence of depression is 5\% higher among SC and 10\% higher among Muslims than higher caste Hindu. 7\% higher anxiety for SC and 10\% higher anxiety for Muslim. 

    \item \cite{ramaiah2015health} the under-five mortality rate (U5MR) was 14 deaths higher among Scheduled Castes (88 vs. 74 nationally), while infant mortality among SCs and STs was 13.1 and 12.3 deaths higher, respectively, compared to the general category (NFHS-3). 
    
    \item \cite{let2024prevalence}: the prevalence of anemia among women of reproductive age was 66.9\% among Adivasi women compared to a national average of 54\%.

\end{itemize}